% Part: incompleteness
% Chapter: introduction
% Section: definitions

\documentclass[../../../include/open-logic-section]{subfiles}

\begin{document}

\olfileid{inc}{int}{def}

\olsection{నిర్వచనలు}

గణితాన్ని సాంకేతికీకరించి, అటువంటి సాంకేతికీకరణ అవైరుధ్యమనీ సంపూర్ణమనీ చూపాలనే
హిల్బర్ట్ కార్యక్రమాన్ని అమలు చేయాలంటే, ముందుగా ఒక భాషనూ తార్కిక చట్రాన్నీ
స్వీకృతాల వ్యవస్థనూ ఎంచుకోవాలి. మన అవసరాల కోసం గణితాన్ని ఒక మొదటిస్థాయి భాషలో
సాంకేతికీకరించవచ్చని అనుకుందాం. అంటే, మొదటిస్థాయి తర్కపు సంయోజకాలూ పరిమాణకాలతో
కలిసి గణిత వాదనలను వ్యక్తం చేయడానికి వీలు కల్పించే కొన్ని
\tetoken{స్థిరసంకేతాలు}{!!{constant}s}, \tetoken{ప్రమేయ సంకేతాలు}{!!{function}s},
\tetoken{విధేయ సంకేతాలు}{!!{predicate}s} ఉన్నాయని అనుకుందాం. అలాంటి భాష ఉందని
చాలామంది అంగీకరిస్తారు: అందులోని ఏకైక తార్కికేతర సంకేతం~$\in$ అయిన సమితి
సిద్ధాంత భాష. ఇంత సరళమైన భాషకు ఇంత వ్యక్తీకరణ శక్తి ఉందనడం మొదటి చూపులో
అసంభవంగా అనిపిస్తుంది; అయినా దాదాపు గణితమంతటినీ ఈ అత్యంత మితమైన పదజాలంలో
వ్యక్తం చేయవచ్చని స్థాపించడానికి చాలా కృషి జరిగింది. విషయాన్ని సరళంగా ఉంచడానికి,
ప్రస్తుతానికి సహజ సంఖ్యలు~$\Nat$ను మాత్రమే వ్యవహరించే గణిత భాగమైన అంకగణితానికి
మన చర్చను పరిమితం చేద్దాం. అంకగణిత వాస్తవాలను వ్యక్తం చేయడానికి సహజమైన భాష
$\Lang L_A$. $\Lang L_A$లో ఒకే ద్విస్థానిక \tetoken{విధేయ సంకేతం}{!!{predicate}}~$<$,
ఒకే \tetoken{స్థిరసంకేతం}{!!{constant}}~$\Obj 0$, ఒక ఏకస్థానిక
\tetoken{ప్రమేయ సంకేతం}{!!{function}}~$\prime$, రెండు ద్విస్థానిక
\tetoken{ప్రమేయ సంకేతాలు}{!!{function}s}~$+$,~$\times$ ఉంటాయి.

\begin{defn}
\tetoken{వాక్యాల}{!!{sentence}s} సమితి~$\Gamma$ పర్యవసానం కింద సంవృతమైతే,
అంటే $\Gamma = \Setabs{!A}{\Gamma \Entails !A}$ అయితే, దాన్ని
\emph{సిద్ధాంతం} అంటాం.
\end{defn}

సిద్ధాంతాలను నిర్దేశించడానికి రెండు సులభ మార్గాలున్నాయి. ఒకటి, ఏదో
\tetoken{నిర్మాణంలో}{!!{structure}} సత్యమైన \tetoken{వాక్యాల}{!!{sentence}s}
సమితిగా నిర్దేశించడం. ఉదాహరణకు, \tetoken{వ్యక్తి క్షేత్రం}{!!{domain}}~$\Nat$
అయి, తార్కికేతర సంకేతాలన్నీ మనం ఆశించినట్లే అర్థనిర్దేశం పొందే $\Lang L_A$
\tetoken{నిర్మాణాన్ని}{!!{structure}} పరిగణించండి.

\begin{defn}\ollabel{def:standard-model}
\emph{అంకగణితపు ప్రామాణిక నమూనా} కింది విధంగా నిర్వచించిన
\tetoken{నిర్మాణం}{!!{structure}}~$\Struct N$:
\begin{enumerate}
\item $\Domain N = \Nat$
\item $\Assign{\Obj 0}{N} = 0$
\item అన్ని $n \in \Nat$లకు $\Assign{\Obj \prime}{N}(n) = n + 1$
\item అన్ని $n, m \in \Nat$లకు $\Assign{\Obj +}{N}(n, m) = n + m$
\item అన్ని $n, m \in \Nat$లకు $\Assign{\Obj \times}{N}(n, m) = n\cdot m$
\item $\Assign{\Obj <}{N} = \Setabs{\tuple{n, m}}{n \in \Nat, m \in
  \Nat, n < m}$
\end{enumerate}
\end{defn}

`$\times$', `$\cdot$'ల మధ్య తేడాను గమనించండి: $\times$ అంకగణిత భాషలోని
సంకేతం. అది గుణకారాన్ని గుర్తుచేయాలని ఎంచుకున్నాం గానీ, $\times$ గుణకార క్రియ
కాదు; అది ఒక ద్విస్థానిక ప్రమేయ సంకేతం (అధికారికంగా, $\Obj f^2_1$). దీనికి
విరుద్ధంగా, $\cdot$ సాధారణ గుణకార ప్రమేయమే. $n \cdot m$ వంటిది కనిపించినప్పుడు
సంఖ్యలు $n$,~$m$ల లబ్ధాన్ని ఉద్దేశిస్తాం; $x \times y$ వంటిది కనిపించినప్పుడు
అంకగణిత భాషలోని ఒక పదం గురించి మాట్లాడుతున్నాం. ప్రామాణిక నమూనాలో ప్రమేయ
సంకేతం~$\times$కు సహజ సంఖ్యలపై ప్రమేయం~$\cdot$ను అర్థనిర్దేశం చేస్తాం. కూడికకు,
అంకగణిత భాషలోని ప్రమేయ సంకేతంగానూ సహజ సంఖ్యలపై కూడిక ప్రమేయంగానూ $+$నే
వాడతాం. ఇక్కడ ఉద్దేశిత అర్థాన్ని సందర్భం నుండి నిర్ణయించాలి.

\begin{defn}
\emph{నిజ అంకగణిత సిద్ధాంతం} అంటే అంకగణితపు ప్రామాణిక నమూనాలో సంతృప్తమయ్యే
\tetoken{వాక్యాల}{!!{sentence}s} సమితి, అంటే
\[
\Th{TA} = \Setabs{!A}{\Sat{N}{!A}}.
\]
\end{defn}

$\Th{TA}$ ఒక సిద్ధాంతం. ఎందుకంటే $\Th{TA} \Entails !A$ అయినప్పుడల్లా,
$\Th{TA}$ను సంతృప్తిపరిచే ప్రతి \tetoken{నిర్మాణంలో}{!!{structure}} $!A$
సంతృప్తమవుతుంది. $\Sat{N}{\Th{TA}}$ కనుక $\Sat{N}{!A}$; అందువల్ల
$!A \in \Th{TA}$.

సిద్ధాంతం~$\Gamma$ను నిర్దేశించే రెండవ మార్గం, ఏదో వాక్య సమితి~$\Gamma_0$
పర్యవసానంగా వచ్చే \tetoken{వాక్యాల}{!!{sentence}s} సమితిగా ఇవ్వడం. అప్పుడు
$\Gamma$, పర్యవసానం కింద $\Gamma_0$ యొక్క ``సంవృతి.'' $\Gamma_0$ను స్పష్టంగా
నిర్దేశించినప్పుడు మాత్రమే, ఉదాహరణకు $\Gamma_0$లోని
\tetoken{మూలకాలను}{!!{element}s} జాబితా చేసినప్పుడు మాత్రమే, ఈ విధంగా సిద్ధాంతాన్ని
నిర్దేశించడం ఆసక్తికరం. కనీసం $\Gamma_0$ నిర్ణయించదగినదై
ఉండాలి; అంటే ఒక \tetoken{వాక్యం}{!!a{sentence}} $\Gamma_0$లోని మూలకమా కాదా
అని నిర్ణయించే గణనీయ పరీక్ష ఉండాలి. $\Gamma_0$లోని
\tetoken{వాక్యాలను}{!!{sentence}s} $\Gamma$కు \emph{స్వీకృతాలు} అంటాం;
$\Gamma_0$ చేత $\Gamma$ \emph{స్వీకృతీకరించబడింది} అంటాం.

\begin{defn}
కింది షరతు నెరవేరితే, అప్పుడు మాత్రమే సిద్ధాంతం~$\Gamma$ను~$\Gamma_0$
\emph{స్వీకృతీకరిస్తుంది}:
\[
\Gamma = \Setabs{!A}{\Gamma_0 \Entails !A}
\]
\end{defn}

\begin{defn}
కింది వాక్యాలు స్వీకృతీకరించే సిద్ధాంతం~$\Th{Q}$ను ``రాబిన్సన్~$\Th{Q}$''
అంటారు; ఇది అంకగణితానికి చాలా సరళమైన సిద్ధాంతం.
\begin{align*}
& \lforall[x][\lforall[y][(\eq[x'][y'] \lif \eq[x][y])]] \tag{$!Q_1$}\\
& \lforall[x][\eq/[\Obj 0][x']] \tag{$!Q_2$}\\
& \lforall[x][(\eq[x][\Obj 0] \lor \lexists[y][\eq[x][y']])] \tag{$!Q_3$}\\
& \lforall[x][\eq[(x + \Obj 0)][x]] \tag{$!Q_4$}\\
& \lforall[x][\lforall[y][\eq[(x + y')][(x + y)']]] \tag{$!Q_5$}\\
& \lforall[x][\eq[(x \times \Obj 0)][\Obj 0]] \tag{$!Q_6$}\\
& \lforall[x][\lforall[y][\eq[(x \times y')][((x \times y) + x)]]] \tag{$!Q_7$}\\
& \lforall[x][\lforall[y][(x < y \liff \lexists[z][\eq[(z' + x)][y]])]] \tag{$!Q_8$}
\end{align*}
\tetoken{వాక్యాల}{!!{sentence}s} సమితి $\{!Q_1, \dots, !Q_8\}$లోని వాక్యాలే $\Th{Q}$కు
స్వీకృతాలు; అందువల్ల వాటి పర్యవసానంగా వచ్చే అన్ని
\tetoken{వాక్యాలు}{!!{sentence}s} కలిసి $\Th{Q}$ అవుతాయి:
\[
\Th{Q} = \Setabs{!A}{\{!Q_1, \dots, !Q_8\} \Entails !A}.
\]
\sourcecorrection{OLTEINCINT-005}{ఆంగ్ల మూలంలోని ``The set ... are the axioms''లో కర్త--క్రియ సంఖ్యాసమ్మతి లేదు. సమితిలోని ఎనిమిది వాక్యాలే స్వీకృతాలని స్పష్టంగా చెప్పాం.}
\end{defn}

\begin{defn}
$!A(x)$ అనేది స్వేచ్ఛా చరాలు $x$, $y_1$,~\dots,~$y_n$ గల $\Lang L_A$లోని
\tetoken{ఒక సూత్రం}{!!a{formula}} అనుకుందాం. అప్పుడు కింది రూపంలోని ప్రతి
\tetoken{వాక్యం}{!!{sentence}} \emph{ఆగమన పథకానికి} ఒక రూపం:
\[
\lforall[y_1][\dots\lforall[y_n][((!A(\Obj 0, y_1, \dots, y_n) \land
\lforall[x][(!A(x, y_1, \dots, y_n) \lif
!A(x', y_1, \dots, y_n))]) \lif
\lforall[x][!A(x, y_1, \dots, y_n)])]]
\]
\sourcecorrection{OLTEINCINT-006}{ఆంగ్ల మూలంలోని సాధారణ ఆగమన పథక ప్రదర్శన $!A$కు స్వేచ్ఛా పారామితులు $y_1,\dots,y_n$ ఉన్నాయని చెప్పినా, మూడు $!A$ రూపాలన్నింటి నుంచీ వాటిని వదిలేసింది. వెంటనే వచ్చే గుర్తింపు పరీక్షలోని సరైన రూపానికి అనుగుణంగా ఆ పారామితులను పునరుద్ధరించాం.}

\emph{పియానో అంకగణితం}~$\Th{PA}$ అంటే, $\Th{Q}$ స్వీకృతాలకూ ఆగమన పథకపు
అన్ని రూపాలకూ స్వీకృతీకరించబడిన సిద్ధాంతం.
\end{defn}

\begin{explain}
ఆగమన పథకంలోని ప్రతి రూపం~$\Struct{N}$లో సత్యం. \tetoken{సూత్రం}{!!{formula}}~$!A$కు
ఒకే స్వేచ్ఛా \tetoken{చరం}{!!{variable}}~$x$ ఉన్నప్పుడు దీన్ని చూడటం అత్యంత
సులభం. అప్పుడు $!A(x)$, $\Struct{N}$లో~$\Nat$ యొక్క ఉపసమితి~$X_{!A}$ను
నిర్వచిస్తుంది. $s(x)=n$ అయినప్పుడు $\Sat{N}{!A(x)}[s]$ అయ్యే అన్ని
$n \in \Nat$ల సమితి $X_{!A}$. దీనికి అనుగుణమైన ఆగమన పథక రూపం
\[
((!A(\Obj 0) \land \lforall[x][(!A(x) \lif !A(x'))]) \lif
  \lforall[x][!A(x)]).
\]
దాని పూర్వాంగం~$\Struct{N}$లో సత్యమైతే, $0 \in X_{!A}$; అలాగే $n \in X_{!A}$
అయినప్పుడల్లా $n+1$ కూడా అందులో ఉంటుంది. $0 \in X_{!A}$ కనుక $1 \in X_{!A}$.
$1 \in X_{!A}$ నుండి $2 \in X_{!A}$ వస్తుంది. ఇలాగే కొనసాగుతుంది. అందువల్ల
ప్రతి $n \in \Nat$కు $n \in X_{!A}$. అంటే $\lforall[x][!A(x)]$,
$\Struct{N}$లో సంతృప్తమవుతుంది.
\end{explain}

$\Th{Q}$, $\Th{PA}$ రెండూ స్వీకృతీకరించిన సిద్ధాంతాలే. అవి ఎంత బలమైనవన్నదే
పెద్ద ప్రశ్న. ఉదాహరణకు, $\Lang L_A$లో వ్యక్తం చేయగల~$\Nat$ గురించిన సత్యాలన్నిటినీ
$\Th{PA}$ నిరూపించగలదా? మరింత నిర్దిష్టంగా, $\Lang L_A$లో చెప్పగల అన్ని
ప్రశ్నలనూ $\Th{PA}$ స్వీకృతాలు పరిష్కరిస్తాయా?

దీన్నే మరో విధంగా, $\Th{PA}=\Th{TA}$ అవుతుందా అని అడగవచ్చు. $\Th{TA}$
$\Nat$ గురించిన సత్యాలన్నిటినీ స్పష్టంగా నిరూపిస్తుంది (అంటే వాటిని కలిగి ఉంటుంది);
$\Lang L_A$లో చెప్పగల అన్ని ప్రశ్నలనూ పరిష్కరిస్తుంది. ఎందుకంటే $!A$ అనేది
$\Lang L_A$లోని \tetoken{ఒక వాక్యం}{!!a{sentence}} అయితే, $\Sat{N}{!A}$ గానీ
$\Sat{N}{\lnot !A}$ గానీ అవుతుంది; అందువల్ల $\Th{TA}\Entails !A$ గానీ
$\Th{TA}\Entails \lnot !A$ గానీ అవుతుంది. అటువంటి సిద్ధాంతాన్ని
\emph{\tetoken{సంపూర్ణం}{!!{complete}}} అంటాం.

\begin{defn}
సిద్ధాంతం~$\Gamma$ భాషలోని ప్రతి \tetoken{వాక్యం}{!!{sentence}}~$!A$కు,
$\Gamma \Entails !A$ గానీ $\Gamma \Entails \lnot !A$ గానీ అయితే, అప్పుడు
మాత్రమే అది \emph{\tetoken{సంపూర్ణమైనది}{!!{complete}}}.
\end{defn}

\begin{explain}
సంపూర్ణతా సిద్ధాంతం ప్రకారం, $\Gamma \Entails !A$ అయితేనే $\Gamma \Proves !A$.
కాబట్టి $\Gamma$ భాషలోని ప్రతి \tetoken{వాక్యం}{!!{sentence}}~$!A$కు,
$\Gamma \Proves !A$ గానీ $\Gamma \Proves \lnot !A$ గానీ అయితేనే అది
సంపూర్ణం.
\end{explain}

మనకు ఎదురయ్యే మరో ప్రశ్న: ఒక \tetoken{వాక్యం}{!!a{sentence}} $\Th{TA}$లో,
$\Th{PA}$లో, లేదా కనీసం $\Th{Q}$లో ఉందో పరీక్షించడానికి గణనాత్మక ప్రక్రియ
ఉందా? ఒక సమితి (ఉదాహరణకు, \tetoken{వాక్యాల}{!!{sentence}s} సమితి) ఎప్పుడు
\tetoken{నిర్ణయించదగినదో}{!!{decidable}} నిర్వచించి ఈ ప్రశ్నను మరింత ఖచ్చితంగా
చెప్పవచ్చు.

\begin{defn}
ఇన్‌పుట్~$x$కు, $x \in X$ అయితే~$1$, లేకపోతే~$0$ను ఇచ్చే గణనాత్మక ప్రక్రియ
ఉంటే, అప్పుడు మాత్రమే సమితి~$X$ \emph{\tetoken{నిర్ణయించదగినది}{!!{decidable}}}.
\end{defn}

కాబట్టి మన ప్రశ్న: $\Th{TA}$ ($\Th{PA}$, $\Th{Q}$)
\tetoken{నిర్ణయించదగినదా}{!!{decidable}}?

ఈ ప్రశ్నలన్నిటికీ సమాధానం: కాదు. ఈ సిద్ధాంతాల్లో ఏదీ నిర్ణయించదగినది కాదు.
అయితే ఈ ఘటన ఈ నిర్దిష్ట సిద్ధాంతాలకే ప్రత్యేకం కాదు. నిజానికి కొన్ని షరతులు
నెరవేర్చే \emph{ప్రతి} సిద్ధాంతానికీ ఇవే ఫలితాలు వర్తిస్తాయి. $\Th{Q}$,
$\Th{PA}$ నెరవేర్చే అటువంటి షరతుల్లో ఒకటి, వాటిని
\tetoken{ఒక నిర్ణయించదగిన}{!!a{decidable}} స్వీకృతాల సమితి స్వీకృతీకరించడం.

\begin{defn}
ఒక సిద్ధాంతాన్ని \tetoken{ఒక నిర్ణయించదగిన}{!!a{decidable}} స్వీకృతాల సమితి
స్వీకృతీకరిస్తే, ఆ సిద్ధాంతం \emph{\tetoken{స్వీకృతీకరించదగినది}{!!{axiomatizable}}}.
\end{defn}

\begin{ex}
ఏ పరిమిత \tetoken{వాక్యాల}{!!{sentence}s} సమితి అయినా
\tetoken{నిర్ణయించదగినది}{!!{decidable}} కనుక, పరిమిత వాక్యాల సమితి
స్వీకృతీకరించే ప్రతి సిద్ధాంతం \tetoken{స్వీకృతీకరించదగినది}{!!{axiomatizable}}.
అందువల్ల, ఉదాహరణకు, $\Th{Q}$ \tetoken{స్వీకృతీకరించదగినది}{!!{axiomatizable}}.

$\Th{PA}$ వంటి పథకరూపంగా స్వీకృతీకరించిన సిద్ధాంతాలు కూడా
\tetoken{స్వీకృతీకరించదగినవే}{!!{axiomatizable}}. $!B$ అనేది $\Th{PA}$
స్వీకృతాల్లో ఒకటో కాదో పరీక్షించడానికి---అంటే $!B$ $\Th{PA}$ స్వీకృతమైతే
$\Char{X}(!B)=1$, లేకపోతే $=0$ అయ్యే ప్రమేయం~$\Char{X}$ను గణించడానికి---ఇలా
చేయవచ్చు: ముందుగా $!B$, $\Th{Q}$ స్వీకృతాల్లో ఒకటో కాదో పరీక్షించాలి. అయితే
సమాధానం ``అవును,'' $\Char{X}(!B)=1$. కాకపోతే, అది ఆగమన పథక రూపమో కాదో
పరీక్షించాలి. దీన్ని క్రమబద్ధంగా చేయవచ్చు; ఈ సందర్భంలో కింది విధానం అత్యంత
స్పష్టంగా ఉండవచ్చు. ఆగమన పథకంలోని ప్రతి రూపం కొన్ని సార్వత్రిక పరిమాణకాలతో
మొదలై, తరువాత ఒక సోపాధికమైన ఉప-\tetoken{సూత్రాన్ని}{!!{formula}} కలిగి ఉంటుంది.
ఆ సోపాధిక ఉత్తరాంగం $\lforall[x][!A(x,y_1,\dots,y_n)]$; ఇక్కడ $x$,
$y_1$,~\dots,~$y_n$ అన్నీ~$!A$ స్వేచ్ఛా చరాలు, $!B$ ముందున్న పరిమాణకాలు
$y_1$,~\dots,~$y_n$ చరాలను బంధిస్తాయి. ఈ~$!A$ను వెలికితీసి, దాని స్వేచ్ఛా
చరాలు ముందున్న సార్వత్రిక పరిమాణకాలూ $\lforall[x]$ బంధించే చరాలకు సరిపోతాయని
తనిఖీ చేసిన తరువాత, సోపాధిక పూర్వాంగం కింది రూపానికి సరిపోతుందో పరీక్షిస్తాం:
\[
!A(\Obj 0, y_1, \dots, y_n) \land \lforall[x][(!A(x, y_1, \dots, y_n)
\lif !A(x', y_1, \dots, y_n))]
\]
అది సరిపోతే $!B$ ఆగమన పథక రూపం; సరిపోకపోతే $!B$ ఆగమన పథక రూపం కాదు.
\end{ex}

ఈ ప్రశ్నకు---అలాగే ఏ సిద్ధాంతాలు సంపూర్ణమో లేదా నిర్ణయించదగినవో అనే మరింత
సాధారణ ప్రశ్నకు---సమాధానం ఇవ్వడంలో కింది నిర్వచనం కూడా ఉపయోగపడుతుంది. సమితి~$X$
ఖాళీగా ఉంటే లేదా \tetoken{ఒక సంగ్రస్త}{!!a{surjective}} ప్రమేయం~$f\colon\Nat
\to X$ ఉంటే, అప్పుడు మాత్రమే అది \tetoken{లెక్కించదగినదని}{!!{enumerable}}
గుర్తు చేసుకోండి. అటువంటి ప్రమేయాన్ని $X$ యొక్క లెక్కింపు అంటాం.

\begin{defn}
సమితి~$X$ ఖాళీగా ఉంటే లేదా దానికి గణనీయ లెక్కింపు ఉంటే, అప్పుడు మాత్రమే దాన్ని
\emph{\tetoken{గణనీయంగా లెక్కించదగిన}{!!{computably enumerable}}} (సంక్షిప్తంగా
\tetoken{గ.లె.}{!!{c.e.}}) అంటాం.
\end{defn}

\tetoken{స్వీకృతీకరించదగినతతో}{!!{axiomatizability}} పాటు, అసంపూర్ణతా సిద్ధాంతాలు
వర్తించే సిద్ధాంతాలపై మరో షరతు: గణనీయ ప్రమేయాలూ
\tetoken{నిర్ణయించదగిన}{!!{decidable}} సంబంధాల గురించిన మౌలిక వాస్తవాలను
నిరూపించగలిగేంత బలంగా అవి ఉండాలి. ``మౌలిక వాస్తవాలు'' అంటే, గణనీయ ప్రమేయాల
ప్రతి ఆర్గ్యుమెంట్‌కు వాటి విలువ ఏమిటో వ్యక్తం చేసే \tetoken{వాక్యాలు}{!!{sentence}s}.
``తగినంత బలంగా'' అంటే, ఆయా సిద్ధాంతాలు ఈ వాక్యాలను తమ సిద్ధాంతఫలితాలుగా
కలిగి ఉండటం. ఉదాహరణకు, ఆదర్శ గణనీయ ప్రమేయమైన కూడికను పరిగణించండి. ఆర్గ్యుమెంట్లు
$2$, $3$లకు $+$ విలువ~$5$, అంటే $2+3=5$. ఆర్గ్యుమెంట్లు $2$, $3$లకు $+$ విలువ~$5$
అని వ్యక్తం చేసే అంకగణిత భాషా వాక్యం: $(\num{2}+\num{3})=\num{5}$. ఉదాహరణకు,
$\Th{Q}$ ఈ వాక్యాన్ని నిరూపిస్తుంది. మరింత సాధారణంగా, ప్రతి గణనీయ ప్రమేయం
$f(x_1,x_2)$కు $\Lang{L_A}$లో \tetoken{ఒక సూత్రం}{!!a{formula}}~$!A_f(x_1,x_2,y)$
ఉండాలని కోరుకుంటాం; $f(n_1,n_2)=m$ అయినప్పుడల్లా $\Th{Q}\Proves
!A_f(\num{n_1},\num{n_2},\num{m})$ కావాలి. ఈ విధంగా, ఆర్గ్యుమెంట్లు
$n_1$,~$n_2$లకు~$f$ విలువ~$m$ అని $\Th{Q}$ నిరూపిస్తుంది. నిజానికి, ఇంకొంచెం
ఎక్కువనే కోరుకుంటాం: $n_1$,~$n_2$లకు~$f$ విలువ మరే సంఖ్యా కాదని కూడా అది
నిరూపించాలి. \tetoken{నిర్ణయించదగిన}{!!{decidable}} సంబంధాలకూ ఇదే వర్తిస్తుంది.
కింది రెండు నిర్వచనాలు దీన్ని ఖచ్చితంగా చెబుతాయి.
\sourcecorrection{OLTEINCINT-007}{ఆంగ్ల మూలంలోని ``the theories in question count these sentences among its theorems''లో బహువచన కర్తకు ఏకవచన సర్వనామం వచ్చింది. ఆయా సిద్ధాంతాలు ఈ వాక్యాలను తమ సిద్ధాంతఫలితాలుగా కలిగి ఉంటాయని ఉద్దేశిత అర్థాన్ని ఇచ్చాం.}

\begin{defn}
$f(n_1,\dots,n_k)=m$ అయినప్పుడల్లా కింది రెండూ నెరవేరితే, అప్పుడు మాత్రమే
\tetoken{సూత్రం}{!!{formula}}~$!A(x_1,\dots,x_k,y)$ ప్రమేయం
$f\colon\Nat^k\to\Nat$కు~$\Gamma$లో
\emph{\tetoken{ప్రాతినిధ్యం వహిస్తుంది}{!!{represents}}}:
\begin{enumerate}
\item $\Gamma \Proves !A(\num{n_1}, \dots, \num{n_k}, \num{m})$, మరియు
\item $\Gamma \Proves \lforall[y](!A(\num{n_1}, \dots, \num{n_k},
y) \lif y = \num{m})$.
\end{enumerate}
\end{defn}

\begin{defn}
కింది రెండూ నెరవేరితే, అప్పుడు మాత్రమే \tetoken{సూత్రం}{!!{formula}}
$!A(x_1,\dots,x_k)$ సంబంధం~$R\subseteq\Nat^k$కు~$\Gamma$లో
\emph{\tetoken{ప్రాతినిధ్యం వహిస్తుంది}{!!{represents}}}:
\begin{enumerate}
\item $R(n_1, \dots, n_k)$ అయినప్పుడల్లా, $\Gamma \Proves !A(\num{n_1},
\dots, \num{n_k})$, మరియు
\item $R(n_1, \dots, n_k)$ కాకపోయినప్పుడల్లా, $\Gamma \Proves \lnot
!A(\num{n_1}, \dots, \num{n_k})$.
\end{enumerate}
\sourcecorrection{OLTEINCINT-008}{ప్రమేయానికి ముందు నిర్వచనం ``in~$\Gamma$'' అని చెప్పినా, సంబంధ నిర్వచనంలోని ఆంగ్ల పరిచయ వాక్యం ఆ నియంత్రక సిద్ధాంతాన్ని వదిలేసింది; రెండు నిబంధనల్లోనూ $\Gamma$ ఉన్నందున దాన్ని పరిచయ వాక్యంలో పునరుద్ధరించాం.}
\end{defn}

ఒక సిద్ధాంతం అన్ని గణనీయ ప్రమేయాలకూ అన్ని
\tetoken{నిర్ణయించదగిన}{!!{decidable}} సంబంధాలకూ
\tetoken{ప్రాతినిధ్యం వహిస్తే}{!!{represents}}, అసంపూర్ణతా సిద్ధాంతాలు వర్తించడానికి
అది ``తగినంత బలమైనది.'' $\Th{Q}$, దాని విస్తరణలు ఈ షరతును నెరవేరుస్తాయి;
కానీ దీన్ని స్థాపించడానికి కొంత సమయం పడుతుంది. $\Th{Q}$ ఎలాంటి విషయాలను
నిరూపించగలదో తెలిపే సామాన్యేతర వాస్తవం ఇది; అన్ని వాస్తవాలనూ నిరూపించడానికి
$\Th{Q}$ వద్ద కొద్ది స్వీకృతాలే ఉన్నాయి కనుక దీన్ని చూపడం కష్టం. అయినా $\Th{Q}$
చాలా బలహీనమైన సిద్ధాంతం. కాబట్టి $\Th{Q}$ అన్ని గణనీయ ప్రమేయాలకు ప్రాతినిధ్యం
వహిస్తుందని నిరూపించడం కష్టమైనా, చాలా ఆసక్తికర సిద్ధాంతాలు~$\Th{Q}$కంటే బలమైనవి,
అంటే $\Th{Q}$ నిరూపించేవాటికంటే ఎక్కువ నిరూపిస్తాయి. $\Th{Q}$ ఏదైనా
నిరూపిస్తే, ప్రతి బలమైన సిద్ధాంతమూ దాన్ని నిరూపిస్తుంది; $\Th{Q}$ అన్ని గణనీయ
ప్రమేయాలకు ప్రాతినిధ్యం వహిస్తుంది కనుక ప్రతి బలమైన సిద్ధాంతమూ అలాగే చేస్తుంది.
అందువల్ల ఎన్నో ఆసక్తికర సిద్ధాంతాలు అసంపూర్ణతా సిద్ధాంతాల ఈ షరతును నెరవేరుస్తాయి.
కాబట్టి మన కఠిన కృషికి ఫలితం ఉంటుంది: అసంపూర్ణతా సిద్ధాంతాలు విస్తృత శ్రేణి
సిద్ధాంతాలకు వర్తిస్తాయని అది చూపుతుంది. ``గణితమంతటినీ'' సాంకేతికీకరించాలని
లక్ష్యంగా పెట్టుకున్న ప్రతి సిద్ధాంతమూ $\Th{Q}$ నిరూపించేవన్నీ తప్పకుండా
నిరూపించాలి; ఎందుకంటే కనీసం ప్రాథమిక గణనల ఫలితాలను అది పట్టుకోగలగాలి. అందువల్ల
``గణితమంతటి'' సిద్ధాంతానికి ప్రతి అభ్యర్థీ అసంపూర్ణతా సిద్ధాంతాలు వర్తించే
సిద్ధాంతమే అవుతుంది.


\end{document}
